\section{固定 BPO 二次结构与全局标量调度}
\label{dyn:scalar}

本节固定一轮 outer 的 prefix 测度 $\mu$、anchor $q$、teacher $t$ 和镜像步长 $\eta>0$。内层优化器学习率记为 $\alpha$。全局标量调度只改变每步的标量步长，保持 AdamW 的一阶与二阶矩规则、betas、epsilon 和更新方向。对 loss 作非线性变换会进入 AdamW 的矩状态，因此另行分析。更新日期为 2026 年 10 月 12 日。

\subsection{精确二次结构及其参数空间边界}

设各分布共享词表且概率为正，$z_q,z_T,z$ 分别为其 logits。定义 $C_q f=f-\E_q f$、$F_q=\operatorname{Diag}(q)-qq^\top$ 与 $z_G=(1-\eta)z_q+\eta z_T$。softmax 的 log 配分函数是与 action 无关的标量，故
\begin{align}
C_q\{\log(p/q)-\eta\log(t/q)\}
&=C_q\{z-z_q-\eta(z_T-z_q)\}\\
&=C_q(z-z_G).
\end{align}
令 $\delta z=z-z_G$。逐 prefix 的平方误差展开为
\begin{align}
\E_q[(C_q\delta z)^2]
&=\sum_a q_a(\delta z_a)^2-\left(\sum_a q_a\delta z_a\right)^2\\
&=\delta z^\top F_q\delta z.
\end{align}
因此
\begin{equation}
L=\frac{1}{2\eta}\E_\mu[\delta z^\top F_q\delta z],\qquad
\nabla_z L=\frac{F_q\delta z}{\eta},\qquad
\nabla_z^2L=\frac{F_q}{\eta}.
\end{equation}
$F_q\mathbf1=0$ 对应公共 logit 偏移。所有曲率与条件数结论均在可辨识的正特征值子空间讨论。有限词表、正 $q$ 下，该子空间的 BPO 最优值为零。$\eta\ne1$ 时零 residual 表示到达 Gibbs target $p_G$，其 teacher reverse KL 一般仍为正；完整 teacher 对齐需要后续 outer 同步。

自由 logits 中还可得到不构造矩阵的统一谱界。对单位向量 $u$，$u^\top F_qu=\operatorname{Var}_q(u_a)$。记其分量范围为 $[l,h]$、均值为 $v$，由 $(u_a-l)(h-u_a)\ge0$ 得
\begin{equation}
\operatorname{Var}_q(u_a)\le(l+h)v-lh-v^2
=\frac{(h-l)^2}{4}-\left(v-\frac{l+h}{2}\right)^2
\le\frac{(h-l)^2}{4}\le\frac12.
\end{equation}
最后一步使用 $(u_i-u_j)^2\le2(u_i^2+u_j^2)\le2\|u\|^2$。所以单 prefix 的 $\lambda_{\max}(F_q/\eta)\le1/(2\eta)$，普通梯度下降的 $0<\alpha<4\eta$ 是一个统一安全区间。若 $u\perp\mathbf1$，则
\begin{equation}
u^\top F_qu=\min_c\sum_aq_a(u_a-c)^2
\ge q_{\min}\min_c\sum_a(u_a-c)^2=q_{\min}\|u\|^2.
\end{equation}
下界随稀有 anchor 动作变小。这个自由 logit 的谱界不包含网络 Jacobian 与 Adam 状态，不能用 $4\eta$ 直接设置生产参数学习率。

若 $z_\theta=z_0+J(\theta-\theta_0)$，则参数问题是精确二次回归，曲率为 $H=\E_\mu[J^\top F_qJ]/\eta$。受限表示下先求二次回归的最优参数 $\theta^*$。一阶最优条件消去交叉项，得到
\begin{equation}
E(\theta):=L(\theta)-L_*=
\frac12(\theta-\theta^*)^\top H(\theta-\theta^*).
\end{equation}
后文的相对误差均针对 $E$。直接使用 $L$ 替代未知的 $E$，会把表示误差底 $L_*$ 误判成可继续消除的误差。

一般神经网络的精确参数 Hessian 为
\begin{equation}
\nabla_\theta^2L=
\frac1\eta\E_\mu\left[J_\theta^\top F_qJ_\theta+
\sum_a(F_q\delta z)_a\nabla_\theta^2z_{\theta,a}\right].
\end{equation}
第一项也随 $J_\theta$ 改变；第二项是 residual 加权的 logit 二阶导数。固定 logit 曲率不等于固定参数曲率。小 residual 或局部线性化可以控制第二项，远离 target 时必须由实际 forward 验证沿方向的二次拟合。

\subsection{绝对下降、相对吸收率与慢方向}

对固定 $H$ 的梯度流，令 $e_i$ 为误差在特征向量上的坐标，正特征值为 $\lambda_i$，有
\begin{equation}
E(\tau)=\frac12\sum_i\lambda_i e_i(0)^2e^{-2\lambda_i\tau}.
\end{equation}
令 $w_i=\lambda_i e_i(0)^2e^{-2\lambda_i\tau}/(2E)$。逐项求导给出
\begin{align}
-\frac{d\log E}{d\tau}&=2\sum_iw_i\lambda_i,\\
\frac{dw_i}{d\tau}&=-2w_i\left(\lambda_i-\sum_jw_j\lambda_j\right),\\
\frac{d}{d\tau}\left[-\frac{d\log E}{d\tau}\right]
&=-4\left(\sum_iw_i\lambda_i^2-\left(\sum_iw_i\lambda_i\right)^2\right)\\
&=-4\operatorname{Var}_w(\lambda)\le0.
\end{align}
单一正特征值下相对吸收率恒定，绝对下降量随 $E$ 减小。多个特征值下快方向先消失，相对吸收率本身也下降。有限维固定谱最终由仍有误差的最小正特征值主导，相对速率趋于正常数；长期趋于零需要越来越小的活跃曲率、连续谱接近零、移动几何或误差底等进一步条件。

在连续时间中对 loss 作递增变换 $\phi(E)$，梯度流满足
\begin{equation}
\dot\theta=-\phi'(E)\nabla E,\qquad
\frac{d\tau}{ds}=\phi'(E).
\end{equation}
只要 $\phi'>0$，路径完全相同，变化是时间参数。理想相对速率设为 $\rho$，则
\begin{equation}
\frac{d\tau}{ds}=\frac{\rho E}{\|\nabla E\|^2}
=\frac{\rho}{2\E_w\lambda}.
\end{equation}
该式需要梯度信息，不能从单个 loss 数值识别 $\E_w\lambda$。一维反例取 $E_h(x)=hx^2/2$，对任意给定 $E_0$ 令 $x=\sqrt{2E_0/h}$，则不同 $h$ 给出相同 $E_0$，而稳定区间为 $0<\alpha<2/h$。任意仅由当前 $E_0$ 指定的正步长都可被足够大的 $h$ 破坏。

\subsection{四类后期变慢的诊断}

固定监督消除了新状态供给这一项，却保留以下四类来源。
\begin{enumerate}
\item KL 与 logit 几何：reverse KL 的梯度为 $F_p(z-z_T)$，一般位置的 Hessian 还包含 $(\partial F_p/\partial z)(z-z_T)$。二分类中 teacher 为均匀分布、student 某一动作概率为 $1-\varepsilon$ 时，logit 梯度量级为 $\varepsilon\log(1/\varepsilon)$，即使 KL 接近 $\log2$，梯度也趋于零。冻结 BPO 把 $F_p$ 换成 $F_q$ 并消去该移动项，但 $q$ 的稀有动作仍对应小曲率。它稳定了 logit 几何，没有消除 anchor 的概率失衡。
\item 曲率谱：固定线性模型的快方向消耗后，剩余误差权重转向低特征值。应测量 $-\Delta\log E$，并与绝对下降 $-\Delta E$ 区分。共享参数的 Jacobian 变化进一步改变谱；在固定误差值上无法识别这两者。
\item AdamW 状态：若二阶矩近似冻结为对角预条件器 $P$，有效曲率是 $P^{1/2}HP^{1/2}$。历史大梯度留下的 $v_n$ 可以压低当前方向步长；$\beta_2=0.999$ 的记忆尺度约为 1000 步，固定数据复用不消除该记忆。动量滞后还会导致 $g^\top d\le0$。若梯度历史不变，改变 lr 只改变当前位置的参数步幅；实际位置改变后未来 moments 也会改变，因此 Adam 轨迹按累计 lr 折叠属于实验结果，不能由 SGD 时间重参数化直接推出。
\item 统计与数值误差底：表示受限产生 $L_*$，固定 sampled pool 产生经验偏差，Top-k center 丢失尾部条件 KL，BF16 产生 forward 分辨率限制，随机 minibatch 或 dropout 才增加跨步噪声。这些项对样本数、支持集大小、forward 精度和 lr 的响应不同。应分别测量完整词表误差、经验 residual 和独立 prefix 的误差。
\end{enumerate}

在偏差修正因子已趋于一的参考模型中，即使 $P$ 固定，Adam 的一阶矩也使每个模式成为二阶递推。记该模式预条件系数为 $p_i$、动量为 $m_n$，则
\begin{equation}
\begin{pmatrix}e_{n+1}\\m_{n+1}\end{pmatrix}
=\begin{pmatrix}
1-\alpha p_i(1-\beta_1)\lambda_i&-\alpha p_i\beta_1\\
(1-\beta_1)\lambda_i&\beta_1
\end{pmatrix}
\begin{pmatrix}e_n\\m_n\end{pmatrix}.
\end{equation}
特征多项式为 $u^2-[1+\beta_1-\alpha p_i(1-\beta_1)\lambda_i]u+\beta_1$。记方括号为 $t_i$，Jury 条件为 $1-\beta_1>0$、$1-t_i+\beta_1>0$、$1+t_i+\beta_1>0$。后两者分别为 $\alpha p_i(1-\beta_1)\lambda_i>0$ 与 $2(1+\beta_1)-\alpha p_i(1-\beta_1)\lambda_i>0$，由此给出 $0<\alpha p_i\lambda_i<2(1+\beta_1)/(1-\beta_1)$。该条件描述固定预条件器动量系统的渐近稳定区间，仍容许瞬时误差上升与振荡。完整 AdamW 的 $P_n$ 随时间改变，这个固定状态计算只用于辨认 moments 的作用。

\subsection{离散步的精确方向二次式}

设一次优化器给出固定方向 $d$，参数候选为 $\theta-\alpha d$。对二次目标直接展开：
\begin{align}
E(\theta-\alpha d)
&=\frac12(e-\alpha d)^\top H(e-\alpha d)\\
&=E-\alpha e^\top Hd+\frac{\alpha^2}2d^\top Hd\\
&=E-\alpha A+\frac{\alpha^2}2B,
\quad A=g^\top d,\quad B=d^\top Hd.
\label{dyn:scalar-ray}
\end{align}
$A\le0$ 时正步长不能减少该二次目标。$A>0,B>0$ 时，沿方向的最优步长及最大相对下降为
\begin{equation}
\alpha_*=\frac AB,\qquad
\rho_{\max}=\frac{A^2}{2BE}.
\end{equation}
由 $H$ 半内积的 Cauchy--Schwarz 不等式，$A^2\le(e^\top He)(d^\top Hd)=2EB$，因此 $0\le\rho_{\max}\le1$。实现某个 $0<\rho\le\rho_{\max}$ 的较小正根为
\begin{equation}
\alpha_\rho=\frac{A-\sqrt{A^2-2B\rho E}}B
=\frac{2\rho E}{A+\sqrt{A^2-2B\rho E}}.
\end{equation}
右式避免相近数相减。加上步长上限时，还需 $\alpha_\rho\le\alpha_{\max}$。给定方向的 $\rho_{\max}$ 下降到目标 $\rho$ 以下后，任何全局标量都无法完成该步的指定比例下降。

对梯度下降 $d=g$，有
\begin{equation}
A=2E\E_w\lambda,\qquad B=2E\E_w\lambda^2,\qquad
\rho_{\max}=\frac{(\E_w\lambda)^2}{\E_w\lambda^2}.
\end{equation}
小步近似忽略 $\alpha^2B/2$，得到 $\alpha\approx\rho E/A$。这就是方向版 Polyak 比例步长的来源；它与最优值已知的经典 Polyak 步长相关，但动量方向下 $g^\top d$ 可以非正，故不能直接套用 SGD 的结论。谱支撑位于 $[m,M]$ 时，精确最速下降的最坏情形相对收缩界为
\begin{equation}
E_{n+1}\le\left(\frac{M-m}{M+m}\right)^2E_n.
\end{equation}
为验证此界，记 $a=\E_w\lambda$、$b=\E_w\lambda^2$。由 $(\lambda-m)(\lambda-M)\le0$ 得 $b\le(m+M)a-mM$，所以
\begin{equation}
\rho_{\max}=\frac{a^2}{b}\ge\frac{a^2}{(m+M)a-mM}.
\end{equation}
右侧函数对 $a\in[m,M]$ 的导数为 $a[(m+M)a-2mM]/[(m+M)a-mM]^2$，最小值位于 $a=2mM/(m+M)$，等于 $4mM/(m+M)^2$。代入 $E^+=E(1-\rho_{\max})$ 得到上述收缩界。

该界仍需要 $O(\kappa\log(1/\varepsilon))$ 次迭代，$\kappa=M/m$。标量线搜索改善步长利用率，但没有把宽谱变成各方向同速。

给出同时占据两个特征值的反例。取 $H=\operatorname{Diag}(m,M)$，在误差 loss 中令 $w_m=M/(m+M)$、$w_M=m/(m+M)$。代入得到 $\rho_{\max}=4mM/(m+M)^2$，谱宽增大时该量为 $O(1/\kappa)$。此时再大的标量也不能沿当前梯度方向实现固定的、与 $\kappa$ 无关的相对下降。

理想化的标量调度可以在离散时间真正改变谱滤波。对固定 $H$、普通梯度下降、已知 $0<m\le\lambda_i\le M$，在 $N$ 步使用
\begin{equation}
\alpha_j=\frac{1}{c-d_0\cos((2j-1)\pi/(2N))},\quad
c=\frac{M+m}{2},\quad d_0=\frac{M-m}{2},
\end{equation}
则误差多项式为
\begin{equation}
P_N(\lambda)=\prod_{j=1}^N(1-\alpha_j\lambda)
=\frac{T_N((c-\lambda)/d_0)}{T_N(c/d_0)}.
\end{equation}
对 $m<M$，其根与常数项分别一致，因而两侧是同一 $N$ 次多项式。$m=M$ 时只需单步 $\alpha=1/m$，单独处理。用 $|T_N(x)|\le1$ 对 $x\in[-1,1]$ 及 $T_N(y)=(u^N+u^{-N})/2$、$u=y+\sqrt{y^2-1}$，得到
\begin{equation}
E_N\le4\left(\frac{\sqrt\kappa-1}{\sqrt\kappa+1}\right)^{2N}E_0.
\end{equation}
Chebyshev 步序需要可信的谱端点、固定线性更新及容许中间放大。它说明标量离散步序理论上具有谱加速能力。当前 AdamW 矩状态、非线性 Jacobian、固定 sampled 更新场及 Top-k center 不满足这个多项式证明的条件，故本节的可实现算法不使用该理想步序。

\subsection{只用观测下降率的反馈规则}

在小步、局部方向变化较慢时，上一更新的观测 log 收缩为 $r_n=\log(E_{n-1}/E_n)\approx\alpha_{n-1}A_{n-1}/E_{n-1}$。要求下一步收缩接近 $r_{\rm target}$，得到
\begin{equation}
\alpha_n=\min\left\{\alpha_{\max},\ \alpha_{n-1}\operatorname{clip}\left(\frac{r_{\rm target}}{r_n},\frac12,\frac54\right)\right\}.
\end{equation}
$E_n\ge E_{n-1}$ 时减半；使用 loss 时应减去已独立验证的 floor。这个规则不增加 backward 或 forward，可以在慢方向占比升高时加大 lr。它沿用上一时刻的方向有效下降率，没有估计当前 $B$；方向转变、动量相位改变或高曲率模式重新进入时无法保证稳定。本节把它作为低开销实验臂，与沿当前方向的 forward 保护规则比较。把 running minimum 当作真实最优值会人为放大相对下降率，故实验不采用该做法。

\subsection{loss 变换与 AdamW 学习率的区别}

设原始梯度为 $g_n$，对整个目标作 $\phi$ 变换后得到 $s_ng_n$、$s_n=\phi'(L_n)>0$。AdamW 的状态为
\begin{equation}
m_n=\beta_1m_{n-1}+(1-\beta_1)s_ng_n,\quad
v_n=\beta_2v_{n-1}+(1-\beta_2)s_n^2g_n^2.
\end{equation}
若从初始化开始 $s_n=s$ 恒定，零初始状态下 $m_n'=s m_n$、$v_n'=s^2v_n$。包含偏差修正后，方向满足
\begin{equation}
\frac{\widehat m_n'}{\sqrt{\widehat v_n'}+\epsilon}
=\frac{\widehat m_n}{\sqrt{\widehat v_n}+\epsilon/s}.
\end{equation}
上述恒等式还要求裁剪未引入不同的跨步缩放。$\epsilon$ 可忽略时常数 loss 缩放几乎不改变 Adam 更新。学习率乘 $s$ 则把实际参数步长乘 $s$，二者并不等价。这个严格恒等式要求相同梯度历史；若先缩放引起轨迹分离，后续原始梯度也会改变。

当 $s_n$ 随 residual 改变时，$m_n'$ 和 $v_n'$ 分别含有不同时间权重的 $s_j$ 和 $s_j^2$。不同参数坐标具有不同梯度历史，所得新方向一般不是原始 Adam 方向的全局标量倍数。若严格要求保留原 AdamW 方向，应在 moment 更新之后改变学习率。非线性 loss 变换只能作为允许改变矩状态时的独立实验臂。

$\phi(E)=\log E$ 与 $\phi(E)=\sqrt E$ 分别给出 $s\propto E^{-1}$、$s\propto E^{-1/2}$。梯度下降下它们在后期增强步长，却没有提供曲率保护；Adam 下增强会被二阶矩部分抵消，还会改变历史响应。以 $\sqrt E$ 为例，梯度流满足 $d\sqrt E/ds=-\E_w\lambda/2$；这控制的是 RMS residual 的绝对下降率，不能保证 $E$ 的相对下降率固定。

单模式二次目标还能说明 lr 的后期趋势由 Adam 状态决定。设 $E=\lambda e^2/2$ 且当前方向为标量 $d$，方向最优步长为 $\alpha_*=e/d$。二阶矩主要保留历史尺度时，$d\approx c e$，对应近似常数 lr；二阶矩追上当前梯度且动量已对齐时，$d\approx\operatorname{sign}(e)$，此时
\begin{equation}
\alpha_*=|e|=\sqrt{2E/\lambda},\qquad
\alpha_\rho=(1-\sqrt{1-\rho})\sqrt{2E/\lambda}.
\end{equation}
后一理想化极限要求 lr 随 RMS residual 缩小，趋势与固定梯度预条件器下补偿慢谱的 lr 增长不同。它依赖当前矩状态、模式曲率和方向对齐，不能作为仅凭 $E$ 的通用规则。三点方向搜索直接测量当前射线，避免预设后期 lr 一定增大或一定减小。

统一 AdamW 学习率时，完整方向为
\begin{equation}
d_n=\frac{\widehat m_n}{\sqrt{\widehat v_n}+\epsilon}+\gamma\theta_n,
\end{equation}
其中 $\gamma$ 为解耦 weight decay。改变 lr 同时改变 decay 的累计作用。即使随机梯度等于目标梯度，动量滞后和 decay 也能使 $g^\top d\le0$。本节实验固定 $\gamma=0$，生产部署如采用非零 decay，试探必须包含该项。

对当前 $g_n$ 的 step，只更新 moments 一次。然后固定 $d_n$，通过 forward 比较候选 $\theta_n-\alpha d_n$。多次试探绝不能多次调用 \texttt{optimizer.step}，否则矩状态、偏差修正计数和方向均改变。拒绝时参数保持原值，已更新的 moments 保留；这等价于该步 lr 为零，下一步继续按正常 AdamW 机制消费梯度。

\subsection{从三个函数值恢复步长}

沿固定方向选择探测尺度 $h>0$，记 $Q_0=Q(0)$、$Q_1=Q(h)$、$Q_2=Q(2h)$。由式~\eqref{dyn:scalar-ray}，
\begin{align}
Q_1-Q_0&=-hA+\frac{h^2}2B,\\
Q_2-Q_0&=-2hA+2h^2B,\\
\widehat B&=\frac{Q_2-2Q_1+Q_0}{h^2},\\
\widehat A&=\frac{Q_0-Q_1}{h}+\frac h2\widehat B.
\end{align}
精确二次目标中两估计严格等于 $B,A$，最优候选为 $\widehat A/\widehat B$。共享参数或经验 residual 平方中，这只是局部拟合。实际算法以 \emph{forward 测量的目标值} 决定接受，拟合本身不作为稳定性证明。

保持当前点为候选，将试探 $h$、$2h$ 与受限拟合候选一并比较；若均未下降，再作有次数上限的二分回退。限制 $\alpha\le\alpha_{\max}$ 和下一次 probe 的增长幅度，使真实非线性目标中的搜索保持局部。探测尺度远小于 logits 的数值分辨率时，二阶差分会受舍入支配；需要相同精度下的常数 lr 对照和完整精度复核。只靠不断缩小 probe 会进入零变化区域。

此算法每次外部迭代使用一次 backward，增加两到三次及回退需要的 forward。减少 backward 次数不自动减少总时间，尤其是模型前向、参数快照和候选安装成本较大时。实际部署可在固定监测子集搜索，但此时只保证该子集的实测下降，必须额外评估独立 prefix 和完整误差。

\subsection{完整 center、固定 sampled pool 与 Top-k 更新场}

先区分总体误差、有限 pool 的经验 residual 平方和生产 surrogate。对完整 center，令 $K_S=\KL(q\Vert p)$、$K_T=\KL(q\Vert t)$，有 $R=U+K_S-\eta(r+K_T)$。总体 $\E_q R=0$，故
\begin{align}
\nabla L
&=\frac1\eta\E_q[R(\nabla\log p+\nabla K_S)]\\
&=\frac1\eta\E_q[R\nabla\log p].
\end{align}
有限、固定 action pool 则一般没有逐 prefix 的经验零均值。经验平方目标的梯度为
\begin{equation}
\nabla\widehat Q=\widehat g+
\frac1\eta\operatorname{mean}_j[R_j\nabla K_{S,j}].
\end{equation}
生产 score surrogate 提供 $\widehat g$，遗漏的第二项在总体 action expectation 下消失，在固定 pool 上不必消失。相同 prefix 重复动作时，它等于该 prefix 的 residual 均值乘 $\nabla K_S/\eta$。因此即使完整 center、固定 Jacobian，沿 $\widehat g$ 的经验平方 residual 也未必下降。对 empirical residual 使用 \texttt{stopgrad} 后作平方，只产生其数值，不产生该平方目标的真实梯度。

生产代码采用冻结 anchor Top-$k$ 支持和一个尾部桶。把该粗粒化 KL 记为 $\widehat K$。KL 的链式分解给出
\begin{equation}
K_S=\widehat K_S+q_{\tail}\KL(q(\cdot\mid\tail)\Vert p(\cdot\mid\tail)).
\end{equation}
因此 $\delta_S=\widehat K_S-K_S\le0$，teacher 的 $\delta_T$ 同理。近似 residual 为 $\widehat R=R+b$，其中 $b=\delta_S-\eta\delta_T$。即使对 action 作完整期望，logit 更新场也变为
\begin{equation}
\E_q[\widehat R\nabla_z\log p]/\eta
=F_q(z-z_G)/\eta+b(q-p)/\eta.
\end{equation}
第二项由条件尾部信息丢失产生。$b$ 的符号没有统一限制，因为它是两个非正量之差。记 $\nabla K_S=p-q$，参数空间扰动为 $-b\nabla K_S/\eta$。其可积性要求 $\nabla b$ 与 $\nabla K_S$ 局部平行；一般粗粒化误差依赖尾部分布形状，满足不了该条件。由乘积求导，反对称 Jacobian 部分包含
\begin{equation}
-\frac1\eta\left[\nabla b(\nabla K_S)^\top-
\nabla K_S(\nabla b)^\top\right],
\end{equation}
故近似场一般不能继承同一个固定平方目标的 Hessian 或单调下降性质。

若全局 loss 缩放仅把原更新场乘以有限且严格正的标量 $s(\theta)$，则 $s(\theta)\widehat g(\theta)=0$ 当且仅当 $\widehat g(\theta)=0$。驻点集合保持不变，所以这类缩放不能把有偏更新场的驻点移回完整 BPO 的驻点。把 lr 设为零只会停止当前位置的更新。AdamW 的瞬态还含矩历史；本节实验中 decay 为零。在学习率保持正且不趋零、更新场连续的常数状态极限中，必要条件仍要求原始梯度为零。

固定 sampled pool 产生固定估计偏差与覆盖不足，重复复用不会将其平均掉。全 batch 确定性重复优化中没有新 minibatch 噪声，不能把这类平台称为随机梯度稳态噪声底。重新划分 minibatch、dropout 或再次抽样才额外产生跨步噪声。若局部更新为 $e_{n+1}=(1-\alpha\lambda)e_n-\alpha\xi_n$，独立零均值噪声方差为 $\sigma^2$，则
\begin{equation}
\operatorname{Var}(e)_\infty=\frac{\alpha\sigma^2}{2\lambda-\alpha\lambda^2}.
\end{equation}
它来自方差递推 $V=(1-\alpha\lambda)^2V+\alpha^2\sigma^2$，增加 lr 会提高误差底。Top-k 偏差、表示误差、浮点误差和 teacher target 未到达则属于不同来源，需要分别记录。

\subsection{最小可实现算法与可证伪结论}

最小算法是对未变换的生产 surrogate 做一次 backward、按标准 AdamW 更新 moments，然后沿固定方向对独立计算的 residual 平方做受限 forward 搜索。方向、betas 和二阶矩公式保持原样。完整目标可用时优先监测完整 BPO；仅 sampled residual 可用时应把该值明确标为经验代理，并同时记录独立误差。使用 $Q(0),Q(h),Q(2h)$ 拟合步长，实测下降才接受，拒绝时采用零 lr 并保留本步 moments。

该算法具有条件性而非普适的加速结论：精确二次目标、下降方向及可分辨的 forward 下，一维拟合得到受限方向最优步长；相对吸收率能否稳定还受 $\rho_{\max}$ 和 lr 上限限制。固定谱下小步调度可只改变数值时间；离散大步与 Adam 矩历史则会改变路径，必须用累计 lr、达到阈值的 backward 次数和独立完整误差共同判定。完整 center 的总体理论不能直接证明固定 sampled/Top-k 生产训练的加速。
